% !TEX root = ../../../main.tex

\subsection{特征数据构成}

实时单词识别模型以 MediaPipe 检测得到的结构化关键点作为训练输入。
系统在服务端调用 MediaPipe Hands 获取左右手关键点，调用 MediaPipe Pose 获取上肢姿态关键点，并将连续帧关键点转换为固定长度的时序特征。
关键点特征能够弱化背景、光照和人物外观差异对识别结果的影响，同时降低输入维度，适合移动端训练反馈场景下的实时识别任务。

本文使用左右手 21 点手部关键点，以及身体上肢的肩、肘、腕关键点作为特征来源。
手部关键点用于描述手型结构、手指弯曲状态和手部局部姿态；
上肢姿态关键点用于描述手腕、手肘相对于肩部中心的位置关系。
实时单词识别所使用的关键点选取如图~\ref{fig:chapter06-realtime-feature-166} 所示。

\WhuAutoFigure{0.82\textwidth}{0.42\textheight}{figures/chapter06/realtime-feature-source-keypoints.png}{实时单词识别特征来源关键点示意图}{fig:chapter06-realtime-feature-166}

手部坐标特征来源于 MediaPipe Hands 输出的 hand world landmarks。
该结果为每只手提供 21 个三维关键点，坐标以米为单位，原点位于手部几何中心附近。
系统对 hand world landmarks 进行尺度归一化，使不同手部大小和不同拍摄距离下的坐标数值具有更接近的尺度范围。

姿态辅助特征来源于 MediaPipe Pose 输出的 pose landmarks。
系统从中选取左右肩、左右肘和左右腕的二维归一化图像坐标，其中 $x$ 和 $y$ 坐标分别按图像宽度和高度归一化到 0 到 1 范围内。
通过上肢关键点，模型能够获得手部动作相对于身体中心的大致位置关系。

在特征组织方式上，系统将一个实时单词识别样本表示为固定长度的时序窗口。
当前模型以连续 30 帧有效图像作为一个输入样本，每帧对应 166 维特征，因此单个样本的输入形状如公式~\ref{eq:chapter06-realtime-sample-shape} 所示。

\begin{equation}
X \in \mathbb{R}^{30\times166}
\label{eq:chapter06-realtime-sample-shape}
\end{equation}

其中，$X$ 表示单个实时单词识别样本，30 表示一个样本包含 30 帧连续有效图像，166 表示每一帧的特征维度。
对于第 $t$ 帧图像，其单帧特征 $f_t$ 由左手特征、右手特征和姿态辅助特征三部分拼接得到，如公式~\ref{eq:chapter06-realtime-frame-feature} 所示。

\begin{equation}
f_t = [H^L_{78},\ H^R_{78},\ P_{10}]
\label{eq:chapter06-realtime-frame-feature}
\end{equation}

表~\ref{tab:chapter06-realtime-feature-composition} 给出了单帧 166 维特征的组成方式。
左右手特征采用相同的维度结构，均由尺度归一化后的手部坐标特征以及手指骨骼角度特征组成；
姿态辅助特征由上肢关键点相对肩中心的位置特征和肘部角度特征组成。

\begin{longtable}{L{0.15\textwidth} L{0.08\textwidth} L{0.25\textwidth} L{0.27\textwidth} L{0.15\textwidth}}
  \caption{单帧 166 维特征构成表}
  \label{tab:chapter06-realtime-feature-composition}\\
  \toprule
  特征部分 & 维度 & 数据来源 & 计算方式 & 作用 \\
  \midrule
  左手坐标特征 & 63 & MediaPipe Hands 的左手 hand world landmarks & 以手腕到中指掌指关节点的距离作为尺度因子，对 21 个三维关键点归一化后展开 & 描述左手空间形态 \\
  左手角度特征 & 15 & 左手 hand world landmarks 的骨骼连接关系 & 基于相邻骨骼向量计算夹角余弦 & 描述左手手指弯曲状态 \\
  右手坐标特征 & 63 & MediaPipe Hands 的右手 hand world landmarks & 以手腕到中指掌指关节点的距离作为尺度因子，对 21 个三维关键点归一化后展开 & 描述右手空间形态 \\
  右手角度特征 & 15 & 右手 hand world landmarks 的骨骼连接关系 & 基于相邻骨骼向量计算夹角余弦 & 描述右手手指弯曲状态 \\
  姿态位置特征 & 8 & MediaPipe Pose 的左右肩、左右肘、左右腕二维坐标 & 以肩中心为参考点进行相对化，并按肩宽进行尺度归一化 & 描述手部相对身体的位置 \\
  姿态角度特征 & 2 & MediaPipe Pose 的左右肩、肘、腕二维坐标 & 计算左右肘部夹角余弦 & 描述上肢弯曲状态 \\
  合计 & 166 & -- & 左右手特征与姿态辅助特征拼接 & 单帧输入特征 \\
  \bottomrule
\end{longtable}

对于单只手，设第 $i$ 个 hand world landmark 为 $h_i=(x_i,y_i,z_i)$，手腕关键点为 $h_0$，中指掌指关节点为 $h_9$。
系统使用手腕点到中指掌指关节点之间的欧氏距离作为手部尺度因子，如公式~\ref{eq:chapter06-hand-scale-factor} 所示。

\begin{equation}
s_h = \lVert h_9 - h_0 \rVert_2
\label{eq:chapter06-hand-scale-factor}
\end{equation}

其中，$s_h$ 表示手部尺度因子。
若尺度因子过小，系统使用极小值 $\varepsilon$ 进行保护，避免除零问题。
归一化后的手部关键点如公式~\ref{eq:chapter06-hand-landmark-normalization} 所示。

\begin{equation}
\hat{h}_i = \frac{h_i}{\max(s_h,\varepsilon)},\quad i=0,1,\ldots,20
\label{eq:chapter06-hand-landmark-normalization}
\end{equation}

其中，$\hat{h}_i$ 表示尺度归一化后的手部关键点。
归一化后的 21 个三维关键点按固定顺序展开，形成 63 维单手坐标特征。

在坐标特征基础上，系统根据 MediaPipe Hands 的手部拓扑结构构造骨骼向量，并对每根手指上的相邻骨骼向量计算夹角余弦。
对于三个连续关键点 $a$、$b$、$c$，以 $b$ 为夹角顶点。
令 $v_1=p_a-p_b$，$v_2=p_c-p_b$，则相邻骨骼向量之间的夹角余弦如公式~\ref{eq:chapter06-finger-angle-cosine} 所示。

\begin{equation}
\cos\theta = \frac{v_1 \cdot v_2}{\lVert v_1\rVert \cdot \lVert v_2\rVert + \varepsilon}
\label{eq:chapter06-finger-angle-cosine}
\end{equation}

每根手指沿骨骼链计算若干相邻骨骼夹角，五根手指共得到 15 维角度特征。
单只手由 63 维坐标特征和 15 维角度特征共同组成 78 维手部特征。

姿态辅助特征用于描述手部动作相对于上半身的位置关系。
系统从 MediaPipe Pose 的 pose landmarks 中选取左右肩、左右肘和左右腕等上肢关键点，先根据左右肩关键点计算肩中心，再将左右手腕、左右手肘转换为相对于肩中心的二维坐标，并按肩宽进行尺度归一化。
设左肩和右肩关键点分别为 $S_l$ 和 $S_r$，肩中心 $C_s$ 的计算方式如公式~\ref{eq:chapter06-shoulder-center} 所示。

\begin{equation}
C_s = \frac{S_l + S_r}{2}
\label{eq:chapter06-shoulder-center}
\end{equation}

身体尺度因子 $s_p$ 使用左右肩之间的欧氏距离表示，如公式~\ref{eq:chapter06-body-scale-factor} 所示。

\begin{equation}
s_p = \lVert S_r - S_l \rVert_2
\label{eq:chapter06-body-scale-factor}
\end{equation}

设上肢关键点集合为 $P_i\in\{W_l,W_r,E_l,E_r\}$，其中 $W_l$、$W_r$、$E_l$、$E_r$ 分别表示左腕、右腕、左肘和右肘关键点，则相对化和尺度归一化后的姿态位置特征如公式~\ref{eq:chapter06-pose-relative-normalization} 所示。

\begin{equation}
\hat{P}_i = \frac{P_i - C_s}{\max(s_p,\varepsilon)},\quad P_i\in\{W_l,W_r,E_l,E_r\}
\label{eq:chapter06-pose-relative-normalization}
\end{equation}

本文取上述 4 个上肢关键点的二维相对坐标，因此姿态位置特征共 8 维。
系统还计算左右肘部夹角，用于描述上肢弯曲状态。
以左侧肘部为例，令 $u_l=S_l-E_l$，$v_l=W_l-E_l$，则左肘夹角余弦如公式~\ref{eq:chapter06-left-elbow-angle-cosine} 所示。

\begin{equation}
\cos\theta_l = \frac{u_l \cdot v_l}{\lVert u_l\rVert \cdot \lVert v_l\rVert + \varepsilon}
\label{eq:chapter06-left-elbow-angle-cosine}
\end{equation}

右侧肘部夹角 $\cos\theta_r$ 的计算方式与左侧相同。
姿态辅助特征由 8 维相对位置特征和 2 维肘部角度特征组成。
最终，单帧 166 维特征可以表示为公式~\ref{eq:chapter06-complete-frame-feature-166}。

\begin{equation}
f_t = [L^{coord}_{63},\ L^{angle}_{15},\ R^{coord}_{63},\ R^{angle}_{15},\ P^{rel}_{8},\ P^{angle}_{2}]
\label{eq:chapter06-complete-frame-feature-166}
\end{equation}

其中，$L^{coord}_{63}$ 和 $R^{coord}_{63}$ 分别表示左、右手尺度归一化后的 63 维坐标特征，$L^{angle}_{15}$ 和 $R^{angle}_{15}$ 分别表示左、右手的 15 维手指骨骼角度特征，$P^{rel}_{8}$ 表示 8 维上肢相对位置特征，$P^{angle}_{2}$ 表示 2 维肘部角度特征。

综上，实时单词识别模型的输入特征由手部局部结构特征和身体上肢姿态辅助特征共同组成。
左手和右手的 78 维特征主要反映经过尺度归一化后的手型结构及手指弯曲状态，10 维姿态辅助特征补充手腕、手肘与肩部中心之间的相对空间关系。
通过将连续 30 帧单帧特征按时间顺序堆叠，系统最终得到形状为 $30\times166$ 的时序输入样本，为后续 1D CNN 模型进行固定词表手势分类提供数据基础。
